\chapter{Conclusion \& Future Work}
\label{sec:conclusion}

\section{Conclusions}

\begin{enumerate}
  \item A causal, wavelength-conditioned light-curve encoder with a
    host-galaxy static branch, trained on simulated LSST-like data and
    fine-tuned on real ZTF data, reaches $\WithStaticMacroF{}$ macro-F1 on
    real, held-out test data (mean of \TierBNSeeds{} seeds), a
    $\SimToRealGapStatic{}$ drop from the same configuration's simulated-data
    score of $\TierLCStaticMacroFFull{}$ (Table~\ref{tab:mainresults}).
  \item Fine-tuning on a comparatively small real labelled set
    (\TierBLabeledTrain{} objects) recovers $\TransferGain{}$ macro-F1 over
    the matching zero-shot, simulation-only model
    (Table~\ref{tab:mainresults}) -- the sim-to-real transfer strategy this
    project depends on is doing real, measurable work.
  \item Host-galaxy static features improve classification on both tiers,
    by $\PhotozGainWith{}$ macro-F1 on simulated data and $\StaticGainTierB{}$
    on real data, and LC+static is the deployed configuration. The
    opposite result reported in an earlier version of this analysis was an
    artifact of a flux-conversion error (Section~\ref{sec:discussion}).
  \item Ib/c F1 on real data, $\WithStaticFIbc{}$, misses the
    \IbcBandLow{}--\IbcBandHigh{} band pre-registered for it. The test split
    holds only \TierBSupportIbc{} Ib/c objects, and the band was derived from
    models with far more Ib/c support; the band is left unchanged and the
    miss is reported.
  \item An anomaly-detection ensemble combining reconstruction error,
    Mahalanobis distance, and an isolation forest reaches
    $\AnomalyTierACombinedThreeRecomputed{}$ ROC-AUC on held-out simulated
    exotic classes and $\AnomalyTierBCombinedThreeRecomputed{}$ on real,
    spectroscopically confirmed non-supernova objects
    (Table~\ref{tab:anomaly}) -- not the best score on either individual
    set, but the only one that stays useful on both.
  \item Classifier confidence is not a usable anomaly signal: it becomes
    \emph{more} confident, not less, on inputs unlike anything the model
    was trained on (Table~\ref{tab:entropybyclass}), and including it in
    the anomaly combination measurably hurt performance on both evaluation
    sets.
  \item The most consequential error in the project passed every test and
    every gate, because the specification, the code and the test all
    shared it. A test of a physical invariant -- that a conversion between
    flux scales must leave a source's magnitude unchanged -- caught it where
    tests of the transcribed formula could not
    (Section~\ref{sec:silentfailures}).
  \item The served pipeline ingests live ZTF and LSST alerts without
    modification, at a median per-alert latency of
    \ServingLatencyPFifty{}\,ms; the LSST alerts are unlabelled, so the
    classifier's accuracy on LSST is untested.
  \item Gate D, the manual inspection of top-ranked anomalies against a
    published real-bogus-artifact anchor, could not be meaningfully
    evaluated: every population available was pre-cleaned of bogus
    candidates before the ranker ever saw it, at two independent pipeline
    stages (Table~\ref{tab:gateoutcomes}).
\end{enumerate}

\section{Future Work}

In order of what would most directly close an open item in this report:
calibrating the classifier's output probabilities, without which no
class-rate estimate from live traffic is meaningful; a fine-tuning
schedule designed for the large brightness shift between simulated and
real objects; closing Gate D against a genuinely unfiltered (or
reject/candidate) alert stream; obtaining labels for LSST alerts so that
accuracy on LSST can be measured; handling LSST's variable-size cutout
stamps; revisiting WHXE now that Tier B fine-tuning has introduced a
genuine population of non-SN-like real objects for its hierarchy to be
tested against; and batching the serving state layer's per-alert database
round trips, since the classifier itself is not the throughput bottleneck.
